\chapter{الكتابة والقراءة}
% full version: chapters/09_acet.tex

لذاكرة الدماغ خاصية مزعجة: الوصلات التي تحفظ الذكريات هي نفسها التي يجري عبرها العمل. لا قرص صلب منفصل. في كل مرة تتذكر الشبكة، تستخدم وصلاتها، وفي كل مرة تتعلم، تغيّرها. فكيف لا تُفسد القديم وهي تكتب الجديد؟ في الحاسوب إشارة منفصلة لهذا: ”السماح بالكتابة“. وفي الدماغ يبدو أن محطة \nt{ACET}، أي الأسيتيل كولين، تؤدي هذا الدور.

\section*{وضعان}

تخيّل شبكة تحفظ بضع صور: أرها نصف صورة فتكمل الباقي. هذه قراءة. والآن أرها صورة جديدة تشبه قديمة. إذا أكملت الشبكة الرسم كعادتها، رأت القديم لا الجديد، وتعلمت الخطأ. كي تكتب الجديد، يجب أن تكف الشبكة مؤقتًا عن الثقة بذاكرتها وأن تنظر إلى العالم.

\pencilfig{9}{\pencilart{9}}{سلك واحد يبدّل الشبكة: أن تكتب الجديد أو تتذكر القديم.}

والأسيتيل كولين يفعل هذا بالضبط، وآثاره الثلاثة كلها مقيسة: يضعف الوصلات الداخلية للشبكة، التي ”تتذكر“ عبرها، ويقوّي الإشارات الآتية من أعضاء الحس، ويسهّل تغيير الوصلات. كثير من الأسيتيل كولين يعني وضع الكتابة: الشبكة تصغي إلى العالم وتتعلم. وقليل منه يعني وضع القراءة: الشبكة تعتمد على الذاكرة وتكمل الرسم.

في نموذجنا لا يوجد مستوى ينجح عنده الأمران معًا. للكتابة يجب خفض الوصلات الداخلية، وللقراءة يُحتاج إليها بكامل قوتها. لذلك لا بد من مبدّل.

\section*{كم نثق بالعينين}

حل المهندسون مسألة مشابهة منذ زمن. حين تسير السفينة في مسارها، يكون لدى الملاح حساب لموضعها المفترض، ولديه أيضًا أرصاد جديدة لكنها غير دقيقة. كم يثق بكل منهما؟ يجيب المرشح الأمثل: ثق بالأرصاد بقدر ما تقل ثقتك بالحساب. وفي الوصفة نفسها أمر ثانٍ: كلما قلت الثقة، وجب إعادة التعلم أسرع. مقبض واحد يتحكم في وزن المدخل وسرعة التعلم معًا. وهذا بالضبط ما يفعله الأسيتيل كولين. ومن هنا الفرضية: إنه يخبر الشبكة بمدى عدم يقين نموذجها الخاص عن العالم.

\pencilfig{124}{%
% optimal blending: the less sure the model, the more weight on the observation (and the faster it relearns)
\foreach \dx/\wm/\we/\ttop/\bbot in {0/2.4pt/0.6pt/{النموذج واثق}/{أسيتيل كولين قليل}, 4.3/0.6pt/2.4pt/{النموذج غير واثق}/{أسيتيل كولين كثير}}{
  \begin{scope}[xshift=\dx cm]
  \draw[penlite=pcViolet, wash=pcViolet, rounded corners=3pt] (0,0.95) rectangle (1.3,1.45); \node[lbl, text=black] at (0.65,1.2) {الذاكرة};
  \draw[penlite=pcBlue, wash=pcBlue, rounded corners=3pt] (0,-0.25) rectangle (1.3,0.25); \node[lbl, text=black] at (0.65,0) {العينان};
  \draw[dotpen=pcAmber, wash=pcAmber] (2.95,0.6) circle (0.55); \node[lbl, text=black] at (2.95,0.6) {التقدير};
  \draw[pcViolet, line width=\wm, line cap=round, -{Stealth[length=5pt]}] (1.35,1.2) -- (2.4,0.8);
  \draw[pcBlue, line width=\we, line cap=round, -{Stealth[length=5pt]}] (1.35,0) -- (2.4,0.4);
  \node[font=\small\itshape, text=pcGray, anchor=south] at (1.55,1.6) {\ttop};
  \node[lbl, anchor=north] at (1.55,-0.4) {\bbot};
  \end{scope}}
}{كم نثق بالعينين. إذا كانت الشبكة واثقة بنموذجها عن العالم، جاء التقدير كله تقريبًا من الذاكرة. وإذا لم تكن واثقة، جاء من الأرصاد، وعندها تعيد الشبكة التعلم أسرع أيضًا. وفق الفرضية، الأسيتيل كولين هو من يدير هذا المقبض.}

\section*{تقسيم العمل}

تذكّر الفصل السابق. حين يتغير العالم فجأة، يلاحظ النورأدرينالين المفاجأة، بحسب الفرضية. أما الأسيتيل كولين فيبقي الشبكة في وضع الكتابة إلى أن تُتعلم البيئة الجديدة. الأول يقول ”شيء ما تغير“، والثاني ”اكتب“.

\section*{إعادة الكتابة الليلية}

في أثناء النوم العميق يكون مستوى الأسيتيل كولين منخفضًا. الشبكة مفصولة عن العالم وفي وضع القراءة. وفي هذا الوقت بالذات ”تُعاد“ فيها أحداث النهار، وهذا مقيس. أما أن ما تُعلّم في النهار يُعاد كتابته عندئذ في الذاكرة الطويلة الأمد، فهي فرضية معقولة، وسنعود إليها في الفصل عن النوم.

\fullversion{9}
